\section*{Bab \ref{ch:g12:deriv} --- Turunan dan Kecembungan}

\begin{solution}{exo:g12:deriv:1}
$f'(x) = 3x^2 - 5$.

Aturan hasil bagi:
$g'(x) = \dfrac{(x^2+1) - x\cdot 2x}{(x^2+1)^2} = \dfrac{1 - x^2}{(x^2+1)^2}$.

Aturan rantai dengan $u = 2x+1$: $h'(x) = 5 \cdot 2 \cdot (2x+1)^4 = 10(2x+1)^4$.

Aturan rantai dengan $u = x^2 + 1$:
$k'(x) = \dfrac{2x}{2\sqrt{x^2+1}} = \dfrac{x}{\sqrt{x^2+1}}$.
\end{solution}

\begin{solution}{exo:g12:deriv:2}
$f'(x) = 2x - 3$, sehingga $f(2) = -1$ dan $f'(2) = 1$: jadi \omterm{def:g12:deriv:tangent}{garis singgungnya}
$y = -1 + (x - 2) = x - 3$. Celahnya adalah
\[
f(x) - (x - 3) = x^2 - 4x + 4 = (x-2)^2 \geq 0 ,
\]
sehingga kurvanya terletak di atas \omterm{def:g12:deriv:tangent}{garis singgungnya} di mana-mana, dan hanya
menyentuhnya di $x = 2$ (seperti diharapkan: $f$ \omterm{def:g12:deriv:convex}{cembung}).
\end{solution}

\begin{solution}{exo:g12:deriv:3}
Keduanya adalah hasil bagi selisih. Dengan $f(x) = (1+x)^{100}$ di $0$:
$f'(x) = 100(1+x)^{99}$, sehingga limitnya $f'(0) = 100$.
Dengan $g(x) = \sqrt{x}$ di $4$: $g'(x) = \frac{1}{2\sqrt{x}}$, sehingga
limitnya $g'(4) = \frac14$.
\end{solution}

\begin{solution}{exo:g12:deriv:4}
\omterm{def:g11:func:function}{Daerah asalnya} $\R \setminus \{1\}$. Limitnya: di $\pm\infty$, $f(x) \sim x \to
\pm\infty$; di $1^\pm$, pembilangnya $\to 4 > 0$ dan penyebutnya $\to 0^\pm$,
sehingga $f \to \pm\infty$; jadi garis $x = 1$ adalah \omterm{def:g12:limcont:asymptote}{asimtot tegak}.

\[
f'(x) = \frac{2x(x-1) - (x^2+3)}{(x-1)^2} = \frac{x^2 - 2x - 3}{(x-1)^2}
= \frac{(x-3)(x+1)}{(x-1)^2}.
\]
Jadi $f' > 0$ pada $\intoo{-\infty}{-1}$ dan $\intoo{3}{+\infty}$, serta
$f' < 0$ pada $\intoo{-1}{1}$ dan $\intoo{1}{3}$. \omterm{def:g11:func:function}{Fungsi} itu naik sampai
\omterm{def:g12:deriv:extremum}{maksimum lokal} $f(-1) = -2$, turun ke $-\infty$, melompat ke $+\infty$ setelah
\omterm{def:g12:limcont:asymptote}{asimtotnya}, turun sampai \omterm{def:g10:functions:extrema}{minimum} lokal $f(3) = 6$, lalu naik.
\end{solution}

\begin{solution}{exo:g12:deriv:5}
Untuk $x \in \intoo{0}{15}$, kotaknya beralas persegi bersisi $30 - 2x$ dan
bertinggi $x$, sehingga
\[
V(x) = x(30 - 2x)^2 .
\]
$V'(x) = (30-2x)^2 + x \cdot 2(30-2x)(-2) = (30-2x)(30 - 2x - 4x)
= (30-2x)(30-6x)$. Pada $\intoo{0}{15}$ berlaku $30 - 2x > 0$, sehingga $V'$
bertanda sama dengan $30 - 6x$: positif sebelum $x = 5$, negatif sesudahnya.
Volumenya \omterm{def:g10:functions:extrema}{maksimum} untuk $x = 5$ cm, yaitu $V(5) = 5 \times 20^2 = 2000\ \text{cm}^3$.
\end{solution}

\begin{solution}{exo:g12:deriv:6}
\emph{1.} $f'(x) = 4x^3 - 12x^2$ dan $f''(x) = 12x^2 - 24x = 12x(x-2)$.
Jadi $f'' > 0$ pada $\intoo{-\infty}{0}$ dan pada $\intoo{2}{+\infty}$ (\omterm{def:g12:deriv:convex}{cembung}),
serta $f'' < 0$ pada $\intoo{0}{2}$ (\omterm{def:g12:deriv:convex}{cekung}). $f''$ berganti tanda di $0$ dan
di $2$: keduanya \omterm{def:g12:deriv:inflection}{titik belok}.

\emph{2.} Di $a = 2$: $f(2) = 16 - 32 + 10 = -6$, $f'(2) = 32 - 48 = -16$,
sehingga \omterm{def:g12:deriv:tangent}{garis singgungnya} $y = -6 - 16(x-2) = -16x + 26$. Celahnya adalah
\[
g(x) = f(x) + 16x - 26 = x^4 - 4x^3 + 16x - 16 .
\]
Karena $g(2) = g'(2) = g''(2) = 0$, maka $(x-2)^3$ membagi $g$; pembagiannya
memberi $g(x) = (x-2)^3(x+2)$. Di dekat $x = 2$ berlaku $x + 2 > 0$, sehingga
$g$ bertanda sama dengan $(x-2)^3$: negatif sebelum $2$, positif sesudahnya.
Jadi kurvanya memotong \omterm{def:g12:deriv:tangent}{garis singgungnya}: $2$ memang sebuah \omterm{def:g12:deriv:inflection}{titik belok}.
\end{solution}

\begin{solution}{exo:g12:deriv:7}
\omterm{def:g11:func:function}{Fungsi} $f(x) = \sqrt{x}$ memenuhi
$f''(x) = -\frac{1}{4}x^{-3/2} < 0$ pada $\intoo{0}{+\infty}$: jadi \omterm{def:g11:func:function}{fungsi} itu
\omterm{def:g12:deriv:convex}{cekung}, sehingga \omterm{def:g10:functions:graph}{grafiknya} terletak \emph{di bawah} setiap \omterm{def:g12:deriv:tangent}{garis singgungnya}.
\omterm{def:g12:deriv:tangent}{Garis singgungnya} di $a = 1$ adalah
\[
y = f(1) + f'(1)(x - 1) = 1 + \tfrac12 (x-1) = \frac{x+1}{2},
\]
sehingga $\sqrt{x} \leq \frac{x+1}{2}$ untuk semua $x > 0$. Kesamaannya berarti
kurvanya menyentuh \omterm{def:g12:deriv:tangent}{garis singgungnya}, dan itu hanya terjadi di titik singgungnya
$x = 1$. (Setara dengan: $\left(\sqrt x - 1\right)^2 \geq 0$.)
\end{solution}

\begin{solution}{exo:g12:deriv:8}
Menurut \cref{thm:g12:deriv:convexity}, \omterm{def:g10:functions:graph}{grafik} \omterm{def:g11:func:function}{fungsi} \omterm{def:g12:deriv:convex}{cembung} yang \omterm{def:g12:deriv:derivative}{dapat
diturunkan} terletak di atas setiap \omterm{def:g12:deriv:tangent}{garis singgungnya}. \omterm{def:g12:deriv:tangent}{Garis singgung} di $a$
mendatar ($f'(a) = 0$), dengan \omterm{def:g10:algebra:equation}{persamaan} $y = f(a)$. Jadi $f(x) \geq f(a)$
untuk semua $x \in \R$: yaitu $f(a)$ adalah \omterm{def:g10:functions:extrema}{minimum} globalnya.
\end{solution}

\begin{solution}{exo:g12:deriv:9}
\emph{Keliling tetap.} Sebuah persegi panjang bersisi $x$ dan $\frac{p}{2} - x$
($0 < x < \frac p2$) berluas $S(x) = x\left(\frac p2 - x\right)$. Lalu
$S'(x) = \frac p2 - 2x$ lenyap di $x = \frac p4$, positif sebelumnya dan
negatif sesudahnya: jadi \omterm{def:g10:functions:extrema}{maksimumnya} di $x = \frac{p}{4}$, tempat kedua sisinya
sama dengan $\frac p4$ --- yaitu sebuah persegi.

\emph{Luas tetap.} Sisi $x$ dan $\frac Ax$ memberi keliling
$P(x) = 2\left(x + \frac Ax\right)$, dengan
$P'(x) = 2\left(1 - \frac{A}{x^2}\right)$, yang negatif untuk $x < \sqrt A$ dan
positif sesudahnya: jadi \omterm{def:g10:functions:extrema}{minimumnya} di $x = \sqrt{A}$, lagi-lagi sebuah persegi.

\emph{Hubungannya.} Kedua pernyataan itu saling dual. Andaikan sebuah persegi
panjang yang bukan persegi meminimumkan kelilingnya pada luas tetap $A$; maka
persegi dengan keliling yang sama akan berluas $> A$ menurut pernyataan
pertamanya, dan menyusutkannya secara sebangun sampai berluas $A$ akan
mengurangi kelilingnya secara tegas --- yang bertentangan dengan
keminimalannya. Jadi masing-masing pernyataan mengakibatkan yang lain.
\end{solution}

\begin{solution}{pb:g12:deriv:1}
\textbf{1.} $30x\left(3x^2 + 1\right)^4$; \quad
$\dfrac{x}{\sqrt{x^2 + 9}}$; \quad
$-\dfrac{2x}{\left(x^2 + 1\right)^2}$.

\textbf{2.} $y(4) = 5$, \omterm{def:g10:reffunc:affine}{gradiennya} $\frac45$: jadi \omterm{def:g12:deriv:tangent}{garis singgungnya}
$y = 5 + \frac45(x - 4)$, yaitu $y = \frac45 x + \frac95$.

\textbf{3.} $f'(x) = \frac{(x^2 + 1) - x \cdot 2x}{(x^2+1)^2}
= \frac{1 - x^2}{(x^2 + 1)^2}$: negatif, positif, negatif
melintasi $-1$ dan $1$: jadi \omterm{def:g10:functions:extrema}{minimumnya} $f(-1) = -\frac12$, \omterm{def:g10:functions:extrema}{maksimumnya}
$f(1) = \frac12$, dengan \omterm{def:g12:limcont:asymptote}{asimtot mendatar} $0$ pada kedua sisinya.

\textbf{4.} $g'(x) = 3x^2 - 6x = 3x(x - 2)$: naik,
turun, naik melintasi $0$ dan $2$; \omterm{def:g12:deriv:extremum}{maksimum lokalnya}
$g(0) = 4$, \omterm{def:g10:functions:extrema}{minimum} lokalnya $g(2) = 0$. $g''(x) = 6x - 6$: \omterm{def:g12:deriv:convex}{cekung}
sebelum $1$, \omterm{def:g12:deriv:convex}{cembung} sesudahnya: jadi \omterm{def:g12:deriv:inflection}{titik beloknya} $(1, 2)$.

\textbf{5.} \omterm{def:g12:deriv:tangent}{Garis singgung} di $1$: \omterm{def:g10:reffunc:affine}{gradiennya} $g'(1) = -3$:
$y = 2 - 3(x - 1) = -3x + 5$. Selisihnya:
$x^3 - 3x^2 + 4 - (-3x + 5) = x^3 - 3x^2 + 3x - 1 =
(x - 1)^3$, yang berganti tanda di $1$: jadi \omterm{def:g12:deriv:tangent}{garis singgungnya}
\emph{memotong} kurvanya --- perilaku khas di sebuah
\omterm{def:g12:deriv:inflection}{titik belok}, tempat kurvanya berpindah sisi.

\textbf{6.} $T(x) = \dfrac{\sqrt{x^2 + 900}}{5} +
\dfrac{\sqrt{(40 - x)^2 + 400}}{2}$. Masuk sebelum $0$ atau
setelah $40$ memperpanjang \emph{kedua} kakinya: jadi sia-sia.

\textbf{7.} $T'(x) = \dfrac{x}{5\sqrt{x^2 + 900}} -
\dfrac{40 - x}{2\sqrt{(40 - x)^2 + 400}}$.

\textbf{8.} Pada segitiga lariannya, sisi sepanjang pantainya
adalah $x$ dan sisi miringnya $\sqrt{x^2 + 900}$: jadi \omterm{def:g11:trigo:cossin}{sinus}
sudutnya terhadap garis tegak lurus persis sama dengan hasil baginya
(sisi depan dibagi sisi miring); serupa dengan itu
$\sin\theta_2 = \frac{40 - x}{\sqrt{(40-x)^2 + 400}}$. Lalu
$T'(x) = 0$ berbunyi
$\frac{\sin\theta_1}{5} = \frac{\sin\theta_2}{2}$ --- yaitu hukum
Snell dengan laju $5$ dan $2$.

\textbf{9.} $T(24) \approx 20.5$ s; $T(40) = 10 + 10 = 20$ s;
$T(0) = 6 + \sqrt{2000}/2 \approx 28.4$ s. Jadi garis lurusnya
bahkan kalah oleh ``berlari sampai ujung, lalu berenang lurus'' --- dan
keduanya kalah oleh jalur yang membias.

\textbf{10.} \omterm{thm:g12:limcont:ivt}{Dikotomi} pada $T'$ (negatif di $24$, positif di
$40$) memusat ke $x^* \approx 34$ m, dengan
$T(x^*) \approx 19.5$ s: sekitar satu detik lebih cepat daripada
garis lurusnya --- yaitu selisih antara sebuah penyelamatan dan sebuah
tragedi.

\textbf{11.} Waktu setiap kakinya merupakan \omterm{def:g11:func:function}{fungsi} \omterm{def:g12:deriv:convex}{cembung} dari $x$
(cabang $\sqrt{\text{bentuk kuadrat}}$), sehingga $T$ \omterm{def:g12:deriv:convex}{cembung}, dan
titik kritis \omterm{def:g11:func:function}{fungsi} \omterm{def:g12:deriv:convex}{cembung} yang \omterm{def:g12:deriv:derivative}{dapat diturunkan} adalah \omterm{def:g10:functions:extrema}{minimum}
globalnya (\cref{exo:g12:deriv:8}): jadi $x^*$ bukan sekadar
stasioner, melainkan tak terkalahkan.

\textbf{12.} Cahaya di dalam air lebih lambat; menurut asas Fermat,
jalur tercepat dari udara ke air membelok di permukaannya tepat menurut
hukum Snell --- sehingga sinar dari pensil yang terendam mencapai mata
dalam keadaan membelok, dan otak, yang menarik garis lurus, melihat
pensilnya patah di garis air.

\textbf{13.} $(x^2)'' = 2 > 0$: jadi \omterm{def:g12:deriv:convex}{cembung}. \omterm{def:g12:deriv:convex}{Kecembungan} di
\omterm{prop:g10:coordgeom:midpoint}{titik tengahnya}: $f\left(\frac{a+b}{2}\right) \leq
\frac{f(a) + f(b)}{2}$, dan itulah pernyataannya. Dengan tangan:
$\frac{a^2 + b^2}{2} - \left(\frac{a+b}{2}\right)^2 =
\frac{(a - b)^2}{4} \geq 0$.

\textbf{14.} Menarik akar kuadratnya (yang naik) pada pertanyaan 13:
$\frac{a + b}{2} \leq \sqrt{\frac{a^2 + b^2}{2}}$. Rantai lengkapnya
pada $2, 8$: harmonik $\frac{2 \cdot 16}{10} = 3.2$; geometri
$\sqrt{16} = 4$; aritmetika $5$; kuadratik
$\sqrt{34} \approx 5.83$: jadi setiap rataan tunduk kepada rataan berikutnya.

\textbf{15.} \omterm{def:g12:deriv:tangent}{Garis singgung} $\frac1x$ di $1$: nilainya $1$, \omterm{def:g10:reffunc:affine}{gradiennya}
$-1$: jadi $y = 2 - x$. Kurva \omterm{def:g12:deriv:convex}{cembung} duduk di atas \omterm{def:g12:deriv:tangent}{garis singgungnya}
(\cref{thm:g12:deriv:convexity}): jadi $\frac1x \geq 2 - x$ untuk semua
$x > 0$, dengan kesamaan tepat di titik sentuhnya $x = 1$.

\textbf{16.} Di \omterm{def:g12:deriv:inflection}{titik beloknya}, cacah kasus \emph{harian}
(yaitu \omterm{def:g12:deriv:derivative}{turunannya}) memuncak: pertumbuhannya berhenti dipercepat lalu mulai
diperlambat --- itulah isyarat matematis pertama bahwa gelombangnya
sedang berbalik, jauh sebelum cacahnya sendiri menurun. Sebelum titik itu,
$f'' > 0$ (setiap harinya lebih buruk daripada hari sebelumnya); sesudahnya,
$f'' < 0$ (masih tumbuh, tetapi lebih lambat).

\textbf{17.} $S'(r) = 4\pi r - \frac{2V}{r^2} = 0$ memberi
$r^3 = \frac{V}{2\pi} = \frac{330}{2\pi}$:
jadi $r \approx 3.74$ cm, dan
$h = \frac{330}{\pi r^2} \approx 7.49$ cm.

\textbf{18.} $h = \frac{V}{\pi r^2} = \frac{2\pi r^3}{\pi r^2}
= 2r$ pada nilai optimumnya: jadi tingginya sama dengan diameternya,
berapa pun volumenya.

\textbf{19.} $S''(r) = 4\pi + \frac{4V}{r^3} > 0$: \omterm{def:g12:deriv:convex}{cembung}, sehingga
jari-jari kritisnya adalah \omterm{def:g10:functions:extrema}{minimum} globalnya. Kaleng yang sesungguhnya
berdiri lebih tinggi demi genggaman, penumpukan, penampakan di rak dan
logam tutup serta alasnya yang lebih tebal --- pengoptimuman selalu
meminimumkan persis apa yang kamu tulis, bukan apa yang kamu maksudkan.

\textbf{20.} Penjaga pantai: modelkan $T(x)$, turunkan, Snell,
\omterm{def:g12:deriv:convex}{kecembungan}, satu detik terhemat. Kaleng: modelkan $S(r)$, turunkan,
$h = 2r$, \omterm{def:g12:deriv:convex}{kecembungan}, logam terhemat. Balok
(\cref{pb:g11:deriv:1}): modelkan $S(w)$, turunkan,
$d = w\sqrt2$, tabel, kekakuan bertambah. Lalu menurut asas
Fermat, cahaya sendiri menjalankan jalur pipa itu di setiap permukaan yang
dijumpainya --- alamlah pengoptimum yang pertama.

\end{solution}
